% Imporditud: tools/import_eksperimendid.py
% Allikas: Eesti füüsikaolümpiaadi eksperimendi ülesannete
%   kogu (Rael Kalda, 2023), ülesanne Ü7, lahendus L7.
\setAuthor{EFO žürii}
\setRound{piirkonnavoor}
\setYear{2000}
\setNumber{P E2}
\setDifficulty{1}
\setTopic{geomeetriline optika}
\setVahendid{Veega täidetud läbipaistev pudel (ilma sildita), paberileht sõnaga AHI}
\setPunktid{11}
\setAllikas{Eksperimendi kogu Ü7, lahendus lk 36}
\setLahendusseis{toores}

\prob{AHI}
Vaadata läbi vett täis pudeli paberile kirjutatud sõna AHI. Kas pudel tuleks asetada piki või risti sõna, et näha sõna IHA. Põhjendage vastust.

\hint

\solu
Pudel peab olema tekstiga risti --- 1 p. Pudel peab olema tekstist kaugel --- 1 p. Vett täis pudel töötab koondava silindrilise läätsena --- 3 p. (kui ainult läätsena, siis 1 p, kui ainult koondava läätsena, siis 2 p). Koondav lääts pöörab kujutise ringi (ümber, vms) --- 1 p. Erinevalt sfäärilisest läätsest, mis pöörab kujutise ringi kõikides suundades, pöörab silindriline lääts kujutist ainult ühe telje ümber, mis ühtib silindri (antud juhul pudeli) teljega --- 2 p. Sõna AHI ongi vertikaaltelje ümber pööratud sõna IHA --- 1 p.
\probend
